\chapter[Sermon 59. Again of the Power of Antichrist, and How the Former Beast Is Worshipped.]{Again of the Power of Antichrist, and How the Former Beast Is Worshipped.}
\label{sermon:059}
\markright{Sermon 59. Again of the Power of Antichrist, and How the Former ...}

And he does all that the first beast could do in his presence. And he causes the earth and those who dwell therein to worship the first beast whose deadly wound was healed.

\index{Antichrist}\index{Augustine, Saint}\index{Arethas of Caesarea}\index{Papacy}\index{Rome}Again he reasons about the power of the second beast, or Antichrist, and of papal authority. He does or executes, he says, the power of the former beast; that is to say, he exercises the same authority that the old Roman Empire exercised. Where he adds, in his presence: Arethas expounds, following immediately after, and even imitating the same. But what power and authority they exercised, I declared before in his place, about the beginning of this chapter. Therefore, just as the Roman Emperors considered all kingdoms and provinces to be theirs, and to belong to them, so the Bishops of Rome claim that all realms are theirs. I give nothing here to affection or hatred. There recently came forth a book printed at Lyons, by Augustine Steuchus, keeper of the Pope’s library, in which he recites from the register of one Gregory (I suppose Gregory VII) all the kingdoms of Europe, Spain, England, France, Denmark, Hungary, etc. The ownership of which belongs to that seat of Rome, the use to the Princes, clients of that see. Often the Popes have attempted to bring the kingdoms of the East under their control also, and that under the pretext of the holy war and recovering the Lord’s sepulchre. And just as the old Romans troubled with continual war the nations that did not acknowledge or obey the old Roman Eagle, so the See of Rome in our time, and in the memory of our forefathers, has put to business and trouble those kingdoms, and nations, and people, who sought to revolt and would not acknowledge those double keys, that is to say, two horns. For who does not know with what cruel wars he vexed in times past the land of Bohemia? Who does not know what Germany and England have suffered in former years? So truly the second beast exercises boldly the tyranny of the old beast. The old beast issued proclamations concerning religion, and the payment of tributes and customs, and in this way impoverished almost all realms, their riches being brought to Rome. And what other thing does that seat do today? What has it done now, to reckon the least, these five hundred years? Who therefore does not see that the second beast exercises most abundantly the power of the first beast? A certain man composed verses in Latin mocking the covetousness and deceits of Rome; and where Rome magnifies herself to be the head of the world, which in Latin is \textit{Caput}, thus he says: If \textit{Caput} comes from \textit{Capio}, which signifies to take, then Rome may well be called so, which does nothing forsake. If you decline \textit{Capio}, \textit{Capis}, and come to the roots, her nets are large and cannot miss, to catch both all and some.

\index{Daniel (Book of)}He adds another point, that this second beast arranges for those who live on Earth to worship the first beast. This, without doubt, we see fulfilled in the Papist kingdom in two ways. First, the Papists have secured such authority and reverence for the Roman Empire, which they call both sacred and holy, that as many as live today, when they hear the name of the Roman Empire spoken, imagine a certain divine thing, brought to them from Heaven. I grant that there have been many truly noble princes, godly, and altogether worthy of praise, within that same Empire: such as Constantine, Augustine, Gratian, Valentinian, Theodosius, and many others. I grant that under these and similar rulers, the Empire was holy, and in truth an empire of Christ. For Christ was acknowledged with true faith; yet we see how the Lord Jesus has nevertheless, as Daniel has done also, called that Empire a beast, doubtless figuratively and for the tyrants. Therefore, we must wisely and justly attribute to each ruler what is his, and not without consideration embrace and reverence that bloody Empire as sacred and holy. And we have declared before, in what manner kingdoms are of God, and how far their works are to be approved that are in kingdoms. And of this there shall be spoken a little later yet more fully.

Secondly, the second beast causes men to worship the first, in this chiefly, that Papistry has brought back the pagan manner, the names only changed. For I told you before, that the first beast was worshipped, in this that simple people received the Roman religion and worshipped idols. The pagans did truly confess the almighty high God, but they joined to him many gods, to whom they submitted elements, diseases, arts, countries, cities, the members and parts of man, and such other like things. Aeolus was god of the winds, and Neptune of the sea, Pluto ruled in the earth, Mars was god of war, Minerva and Apollo of arts, Aesculapius over diseases, Hercules and many more. Venus was lady of love, and the goddess Juno of marriage. Neither was there any member in the body that did not have his god: so had all countries and cities their saving gods, and every house its domestic gods. To them afterward they framed idols, that is tokens and memorials, which might bring those heavenly gods into the memory of the earthly dwellers. They built chapels and churches, they instituted priests, holy days, rites, and ceremonies. These things are found in the books of the gentiles, and in our histories, and also in the writings which have refuted the pagan idolaters. But in the papal kingdom at this day, the names being only changed, who can deny that the same cult, the same religion, indeed very superstition is not renewed? Of these things I have treated at length in my book \textit{De origin erroris}. The Papists teach that the saints in heaven reign with God, and that to them are subject sicknesses, arts, limbs or members, cities, and all things, and must therefore be called upon and worshipped. Saints are expressed and represented by images, to these images are erected altars and churches: briefly, it is done to them that was done to the gods and idols of the pagans. Who therefore understands now that Antichrist has procured that the first beast might be worshipped, that is to wit, might be powerful again, and that the old idolatry and superstitious worship might be renewed and frequented?

And as we have read before, it is said that all who dwell on Earth worship him, those whose names are not written in the book of life from the Lamb. He says this plainly here as well, and he causes the Earth and the inhabitants of the Earth—that is, those who seek and regard only earthly things—to worship the first beast. Not all are polluted with papal idolatry. To this pertains the noble history of Leo the third, Emperor, and Gregory the second, and of other Popes, through whose wickedness idolatry was again brought into the church; of which I wrote at length in my work \textit{On the Origin of Errors}.

Neither is this annexed here without a mystery, whose deadly plague was healed. For he seems to compare together the first and second beast, and to show their likeness. And I told you how many men at the first were kept still in the Roman errors and idolatry, because the gods, by Vespasian’s means, were said to have preserved the commonwealth, which else with civil wars was as it were brought to ruin. Finally, we read in histories that the Empire of Rome has many times received deadly wounds; but yet, immediately, through the wisdom and valor of some noble men, the gods (as they speak) so willing, have been healed again. In that number are Lucius Septimius Severus, Valerius Aurelianus, C. Aurel. Val. Diocletian, and so forth. By whose lucky success, triumphs, and victories restoring the Empire, many have been moved to say, who does not see that Rome shall be eternal, and that the Roman religion is to the gods most acceptable, and that the emperors also and public weal is endowed with a certain deity, and is to be honored? After the same sort, the kingdom of the Pope or Antichrist, having tried diverse chances, has very often escaped out of desperate dangers. Force and policy have afflicted it, and also the religion of Henry III, Emperor, and of his son Henry IV, Frederick I and II vexed the popes. There were also other mighty princes who inflicted mortal wounds to the See of Rome.

Again there were bishops of Rome who with great cunning have healed their wounds. Among these were Gregory VII, Urban II, Paschal II, Calixtus II, Alexander III, Innocent III, Honorius III, Gregory IX, Clement IV and V, Boniface VIII, John XXII, and many others. But was that situation not in greatest danger in times past, when three popes were created at once: one residing in Rome, the second going to Avignon in France, and the third living in Spain? But all three were deposed by the power, diligence, authority, and policy of Emperor Sigismund and the council of Constance. That deadly wound was fairly healed in Martin V. And this felicity, and the restoration of the papal kingdom, persuades many effectively that popery is of God, and the popish religion to be most certain and true: as that which has often been attacked by mighty princes might indeed be shaken, but never yet overthrown. The acclaim of all the Romanists is known: the ship of St. Peter is indeed tossed with storms, but cannot be drowned. But Daniel himself also prophesied that this would come to pass, saying: “He shall prosper and shall do what he will, and shall kill the strong and holy people at his pleasure, and guile shall be directed in his hand.” Those who are offended at this day by the felicity of that chair of pestilence and the beast thereof do not mark these things. Therefore, just as the days of mourning and sudden destruction came upon old Rome and utterly destroyed both the city and empire, even so we shall hear in the 17th and 18th chapters that Babylon shall have her fatal destinies. The Lord Jesus confirm us in the faith of Jesus Christ, and deliver us from the guiles, lucky success, and felicity of that Roman Antichrist. Amen.
